%root = ../thesis-main.tex
\chapter{Designing a Reactive Architecture for Alchemist}\label{design_alchemist}

\section{Alchemist}
Alchemist~\cite{pianini2013COSCSA} is a chemical-oriented, general-purpose
stochastic simulator, extensible and open source developed by the University of
Bologna that runs on the \ac{JVM}, built to tackle the complexity that arises
with \emph{pervasive computing}. This simulator focusses on the necessity to
define advanced mechanisms for modelling autonomous agents' interactions in
highly-dynamic and distributed environments. Thanks to this, Alchemist is
capable to express various scenarios for pervasive, aggregate and
nature-inspired computing.

Alchemist's core is composed of a set of possibly heterogeneous nodes capable
of communicating among themselves through a set of chemical-inspired
interactions and reactions. Such extensible and general-purpose model can lead
to the definition of autonomous behaviours of situated agents, reflecting with
high-fidelity the intrinsic dynamics and complexity that arises in pervasive
contexts. Alchemist's flexibility lets modelers define simulation logic through
a set of behavioural rules on top of the environment, resulting in a powerful
tool for exploring and lead experiments on coordination strategies,
self-adaptation and self-management in multi-agent and pervasive contexts.

\paragraph{Meta-model} The Alchemist meta-model represents the theoretical
foundations which guides simulation modelers. The components of Alchemist are
depicted in \Cref{fig:alchemist-metamodel}:
%
\centerimagesource{figures/alchemist-metamodel.png}{The Alchemist
	meta-model}{alchemist-metamodel}{1}{https://alchemistsimulator.github.io/explanation/metamodel/}
%
\begin{itemize}
	\item \textbf{Molecule} the name associated to a value;
	\item \textbf{Concentration} the value associated to a specific molecule;
	\item \textbf{Node} molecules and reactions' container, situated in an
	      environment;
	\item \textbf{Environment} space abstraction inside Alchemist. It contains
	      all nodes and can perform:
	      \begin{itemize}
		      \item compute nodes' position in space;
		      \item compute the distance between two nodes;
		      \item move nodes inside the environment.
	      \end{itemize}
	\item \textbf{Linking Rule} a function of the current status of the
	      environment. It associates a neighborhood to every node;
	\item \textbf{Neighborhood} entity composed by a node (\emph{center}) and a
	      set of nodes (\emph{neighbors});
	\item \textbf{Reaction} represents any event that may change environment's
	      state. Every reaction is defined by a set of conditions (even none),
	      one or more actions and a \emph{Time Distribution}.
	      \Cref{fig:alchemist-reaction} shows components of a reaction. The
	      frequency at which at which a reaction can happen could depend on a
	      static ``rate'' parameter, the value of each condition or a ``rate
	      equation '' which combines the static rate and condition's propensities
	      to compute an ``instantaneous rate'';
	\item \textbf{Condition} a function on current environment's state yielding
	      a boolean and a real number: if the condition does not hold, the
	      associated reaction will not execute. The real number returned by a
	      condition may influence reaction's speed, depending on the reaction
	      itself and the time distribution;
	\item \textbf{Action} models a change on the environment.
\end{itemize}
%
\centerimagesource{figures/alchemist-reaction.png}{Reaction representation in
	Alchemist}{alchemist-reaction}{1}{https://alchemistsimulator.github.io/explanation/metamodel/}
%
\paragraph{Incarnations} Alchemist's \emph{extensibility} relies on the
abstract representation of molecules and concentrations. This is realized
through the concept of \emph{Incarnations}, which provides concrete type
definitions for molecules and concentrations. An \emph{Incarnation} may also
specify conditions, actions, and occasionally, the environment and the
reactions that act upon it. In short, an \emph{Incarnation} is a concrete
instance of Alchemist's meta-model. Although Alchemist employs classical
chemical terminology, different \emph{Incarnations} can represent vastly
different universes. More specifically, at the time of writing, Alchemist
supports the following \emph{Incarnations}:

\begin{itemize}
	\item \textbf{SAPERE}~\cite{pianiniCISFPSE} was the first stable
	      \emph{Incarnation} for Alchemist, and defines a nature-inspired
	      framework designed to engineer self-organising pervasive service
	      ecosystems. In this approach agents (such as device or users) are
	      represented by Live Semantic Annotations (LSAs), which interact within
	      a distributed spatial substrate governed by the so called ``eco-laws''.
	      Functioning like chemical reactions, these laws enable spontaneous
	      discovery, spatial diffusion and data aggregation without centralised
	      control. Its practical viability is demonstrated through a large-scale
	      real-world deployment at the Vienna City Marathon, where it
	      successfully managed \emph{crowd steering} via public displays and
	      smartphones~\cite{zambonelli2015DPMASNIC}.
	\item \textbf{Protelis}~\cite{pianini2015PROTELIS}: Protelis is a
	      functional language embedded in Java designed for practical
	      \emph{aggregate programming}. This paradigm shifts the focus from
	      individual device behaviour to the computation and coordination of
	      \emph{aggregates of devices}, enabling the decomposition of problems
	      solvable by distributed networks. Protelis relies on the \emph{Field
		      Calculus} to ensure that these aggregate specifications are
	      consistently mapped to local executions, allowing for the construction
	      of resilient, widely reusable distributed components.
	\item \textbf{Biochemistry} is an Alchemist incarnation created to support
	      \emph{biochemical} reactions which occur inside a biological cell or a
	      group of cells that share a common environment.
	\item \textbf{\textsc{ScaFi}} is a Scala toolkit providing an internal
	      domain-specific language, a simulation environment, and runtime support
	      for practical aggregate computing. This framework implements a variant of
	      the Higher-Order Field Calculus, offering a unified platform for
	      simulating and executing distributed collective systems.
\end{itemize}

\subsection{Dependency Management in Alchemist}
The stochastic chemical simulation concepts presented in
\Cref{sec:dependency_management} assume single compartments of interacting
molecules. Alchemist extends this domain by introducing the concept of
``neighborhoods,'' defined as sets of communicating \emph{Nodes}. In this
scenario, nodes can influence one another and move within an environment.
However, since molecules can travel between compartments and nodes can move
relative to one another, maintaining the dependency graph becomes a significant
challenge. In addition to the standard \ac{SSA} dependencies, this abstraction
introduces \emph{spatial dependencies}, transforming the structure into a
\emph{spatial dependency graph} where reaction propensities may be affected
locally, within a neighborhood or globally. To manage this complexity,
Alchemist reactions define specific \emph{input} and \emph{output} contexts.
These contexts specify, respectively, which parts of the environment a reaction
reads to compute its propensity and which parts it modifies.

Alchemist addresses this challenge by implementing the \emph{Next Reaction
Method} (\Cref{ssec:ssa_optimisation}), which explicitly maps reaction
dependencies in a directed graph. Maintaining the consistency of this graph is
a non-trivial task. Specifically, adding a reaction requires verifying its
dependencies against existing system reactions; if dependencies are found, a
new node is added to the graph with corresponding edges. Conversely, deleting a
reaction requires pruning all dependencies where the reaction acts as either
the influencer or the influenced. Notably, any topology change during the
simulation (e.g., neighbors entering or leaving a neighborhood) triggers these
procedures, requiring a dynamic recalculation of dependencies for the affected
components.

\subsection{In-house Solution} % should this be moveod into implementation detail? it's more a design aspect...
The analysis of these standards reveals a significant gap between industrial
reactive tools and the rigorous requirements of a Discrete Event Simulator.
Frameworks like Akka and RxJava prioritize throughput and availability,
design choices that inherently allow for non-deterministic event interleaving.
Conversely, a simulation engine requires strict causal ordering and glitch
freedom to accurately calculate system dynamics; relying on eventual consistency
is insufficient for physical correctness. Furthermore, commercial frameworks
typically couple execution to wall-clock time, whereas our engine necessitates
a \emph{virtual time} scheduler to decouple simulation progression from hardware
execution speed.

To address these limitations, we propose an in-house reactive engine. This
solution adopts the expressive vocabulary of RxJava (Observables, Signals) and
the deterministic rigor of academic models like Lingua Franca~\cite{menard2023LF},
implementing a propagation strategy specifically tailored to the discrete-event
lifecycle.
